% from ega4-21c-fr.tex
\oldpage[IV-4]{318}
$\mathrm H^1(U_0,\mathcal O_{U_0})=0$ en raison de la relation
$\operatorname{prof}(B)\geq3$. Par ailleurs, comme $B$ est factoriel, on a
$\mathrm H^1(U,\mathcal O_{U_0}^*)=\operatorname{Pic}(U_0)=0$; on tire donc bien de
la suite exacte précédente que
$\operatorname{Pic}(U)=\mathrm H^1(U,\mathcal O_U^*)=0$.

(iv) Il existe aussi des anneaux locaux noethériens de dimension $3$ qui sont
intègres, intégralement clos et parafactoriels, mais non factoriels. Soit en effet $B$
un anneau local noethérien, complet, intègre et intégralement clos, de dimension
$\geq2$ et non factoriel (par exemple le complété de l'anneau localisé de l'algèbre
intègre $k[U,V,W]/(W^2-UV)$ en l'idéal image de $(U)+(V)+(W)$). Alors il résulte de
ce que l'on verra plus bas (21.14.2) que l'anneau local $A=B[[T]]$ des séries formelles
en $T$ sur $B$ est parafactoriel, mais il n'est pas factoriel, sans quoi $B$ le serait
en vertu de (21.13.12) ci-dessous.

(v) On peut montrer qu'un anneau local d'intersection complète absolue (19.3.1) de
dimension $\geq4$, est parafactoriel (cf. [41, XI, 3.13 (i)]).

(vi) On a vu (Remarque (ii)) que tout anneau local noethérien factoriel $A$ de
dimension $2$ est parafactoriel. Mais il y a des anneaux locaux noethériens de
dimension $2$ qui sont parafactoriels et non factoriels. On peut montrer que, pour
qu'un anneau local noethérien $A$ de dimension $2$ soit parafactoriel et non
factoriel, il faut et il suffit qu'il satisfasse aux trois conditions suivantes :
\begin{enumerate}[label=\arabic*\textsuperscript{o}]
  \item $A$ est un anneau de Cohen--Macaulay (autrement dit
  $\operatorname{prof}(A)=2$).
  \item $A$ est intègre et si $A'$ est sa clôture intégrale, $A'$ est factoriel et est
  une $A$-algèbre finie.
  \item Soit $\mathfrak J$ le conducteur de $A$ dans $A'$ (annulateur du $A$-module
  $A'/A$, ou encore le plus grand idéal de $A'$ contenu dans $A$); posons
  $B=A/\mathfrak J$, $B'=A'/\mathfrak J$; alors $\dim(B)=1$ (ce qui implique
  $A'\neq A$, autrement dit $A$ n'est pas intégralement clos), et l'application
  canonique $\operatorname{Div}(B)\to\operatorname{Div}(B')$ (21.4.5) est
  \emph{surjective}.
\end{enumerate}
On peut montrer en outre que ces conditions entraînent la propriété suivante :

$4^{\mathrm o}$ L'anneau $B'$ (et \emph{a fortiori} $B\subset B'$) est réduit, et le
morphisme $\operatorname{Spec}(B')\to\operatorname{Spec}(B)$ est bijectif.

Si l'on pose $X=\operatorname{Spec}(A)$, $X'=\operatorname{Spec}(A')$,
$Y=\operatorname{Spec}(B)$, $Y'=\operatorname{Spec}(B')$, de sorte que $Y$
(resp.\ $Y'$) est défini par l'idéal $\mathfrak J$ de $A$ (resp.\ $A'$), le morphisme
structural $f:X'\to X$ est un isomorphisme de $X'-Y'$ sur $X-Y$ (Bourbaki,
\emph{Alg. comm.}, chap.~V, \S1, n\textsuperscript{o}~5, cor.~5 de la prop.~16). On
voit donc en vertu de $4^{\mathrm o}$ que $f$ est un morphisme \emph{bijectif}; en
d'autres termes, $X$ est un préschéma \emph{unibranche} (6.15.1) (et en particulier
$A'$ est un anneau local noethérien); en général $X$ n'est pas géométriquement
unibranche. L'espace $Y$, de dimension $1$, est constitué par le point fermé $x$ de
$X$ et les points maximaux $y_i$ ($1\leq i\leq r$) de $Y$, et l'unique point $y_i'$
de $Y'$ au-dessus de $y_i$ est aussi un point maximal de $Y'$; comme $B$ et $B'$ sont
réduits, on en déduit que l'on a
\[
  \mathfrak m_{y_i}\mathcal O_{X',y_i'}=\mathfrak m_{y_i'}.
\]
Pour $z=y_i$ et $z'=y_i'$, cette relation est aussi vérifiée de façon évidente pour
$z\notin Y$ et $z'$ l'unique point de $f^{-1}(z)$, donc pour tout
$z\in U=X-\{x\}$ et $z'$ l'unique point de $f^{-1}(z)$. En particulier, si l'anneau
$A$ est de caractéristique $0$ (0, 21.1.1) on voit que $U'$ est \emph{non ramifié} sur
$U$ (mais non étale en général).

Nous nous bornerons ici à démontrer que les conditions $1^{\mathrm o}$,
$2^{\mathrm o}$, $3^{\mathrm o}$ ci-dessus sont \emph{suffisantes} pour que $A$ soit
parafactoriel. Or, en vertu de la condition $1^{\mathrm o}$ et de (21.13.8), il suffit
de démontrer que $\operatorname{Pic}(U)=0$. Puisque $A'$ est factoriel en vertu de
$2^{\mathrm o}$, on a $\operatorname{Pic}(U')=0$, et on déduit de (21.8.5, (ii)) une
suite exacte
\[
  1\longrightarrow
  \left(\prod_{i=1}^r
  \mathcal O_{U',y_i'}^*/\mathcal O_{U,y_i}^*\right)
  \big/\operatorname{Im}\Gamma(U',\mathcal O_{U'}^*)
  \longrightarrow\operatorname{Pic}(U)\longrightarrow\operatorname{Pic}(U')
  \longrightarrow0
\]
et il s'agit donc de montrer que le second terme de cette suite est réduit à l'élément
neutre, en prenant pour $S$ l'ensemble des $y_i$ ($1\leq i\leq r$) et en notant
qu'ici $f_*(\mathcal O_{X'})$ est la $\mathcal O_X$-Algèbre $\widetilde{A'}$. Soient
$\mathfrak p_i$, $\mathfrak p_i'$ les idéaux premiers de $A$ et $A'$ correspondant
aux points $y_i$, $y_i'$, et remarquons que $\mathfrak p_i'\cap A=\mathfrak p_i$;
donnons-nous pour chaque $i$ un élément inversible $a_i'/s_i$ de
$A'_{\mathfrak p_i'}$, avec $a_i'\in A'$, $s_i\notin\mathfrak p_i'$; comme nous
n'avons à examiner que les groupes quotients
$(A'_{\mathfrak p_i'})^*/(A_{\mathfrak p_i})^*$, on peut supposer $s_i=1$ pour tout
$i$; soit $b_i'\in B'=A'/\mathfrak J$ l'image canonique de $a_i'$, de sorte que
$b_i'/1$ est un élément non nul de $k(y_i')=\mathcal O_{Y',y_i'}$. L'ensemble
$U\cap Y$ formé par les $y_i$ est un ouvert affine de la forme $D(t)$ dans $Y$, avec
$t\in\mathfrak m$, idéal maximal de $A$; il existe donc un élément inversible
$b'\in B'_t$ dont les $b_i'/1$ sont les images canoniques, et comme $t$ est régulier
dans l'anneau intègre $A'$, on peut supposer $b'$ de la forme $b''/t^m$, où
$b''\in B'$ est inversible. Soit $a''$ un élément de $A'$ dans la classe $b''$, qui
est donc nécessairement inversible; $a'=a''/t^m$ est inversible dans $A'_t$ et pour
tout $i$, on a
\[
  u_i=a''-a_i't^m\in\mathfrak J,
  \qquad
  t^m a_i'=a''-u_i=a''(1-a''^{-1}u_i);
\]
mais $a''^{-1}u_i\in\mathfrak J\subset\mathfrak m$, donc $1-a''^{-1}u_i$ est un
élément inversible de $A$, et les classes de $a'$ et $a_i'$ dans
$(A'_{\mathfrak p_i'})^*/(A_{\mathfrak p_i})^*$ sont les mêmes, ce qui achève la
démonstration.

